% !TEX program = pdflatex
\documentclass[10pt]{article}

\usepackage[utf8]{inputenc}
\usepackage[T1]{fontenc}
\usepackage[english]{babel}
\usepackage{lmodern}
\usepackage{microtype}
\usepackage[margin=0.72in]{geometry}
\usepackage{amsmath,amssymb,amsthm,mathtools}
\usepackage{booktabs}
\usepackage{xcolor}
\usepackage{enumitem}
\usepackage[most]{tcolorbox}
\usepackage{fancyhdr}
\usepackage[hidelinks]{hyperref}

\definecolor{CourseBlue}{HTML}{1A3C68}
\definecolor{CourseGreen}{HTML}{216E39}
\definecolor{CourseGold}{HTML}{B8860B}
\definecolor{CourseRed}{HTML}{9A2C2C}
\definecolor{SoftBlue}{HTML}{F2F6FA}
\definecolor{SoftGreen}{HTML}{F1F8F3}
\definecolor{SoftGold}{HTML}{FFF8E8}
\definecolor{SoftRed}{HTML}{FFF3F3}

\setlength{\parindent}{0pt}
\setlength{\parskip}{0.32em}
\setlist[itemize]{leftmargin=1.5em,itemsep=0.12em,topsep=0.2em}
\setlist[enumerate]{leftmargin=1.7em,itemsep=0.12em,topsep=0.2em}
\allowdisplaybreaks

\pagestyle{fancy}
\fancyhf{}
\lhead{Stochastic Calculus and Option Pricing}
\rhead{Lecture 4 Notes}
\cfoot{\thepage}
\renewcommand{\headrulewidth}{0.3pt}

\newcommand{\E}{\mathbb E}
\newcommand{\Pp}{\mathbb P}
\newcommand{\Q}{\mathbb Q}
\newcommand{\R}{\mathbb R}
\newcommand{\1}{\mathbf 1}
\newcommand{\cF}{\mathcal F}
\newcommand{\dd}{\,\mathrm d}
\newcommand{\Ito}{\mathrm{It\hat{o}}}
\DeclareMathOperator{\Var}{Var}
\DeclareMathOperator{\Cov}{Cov}
\DeclareMathOperator{\Med}{Med}

\theoremstyle{definition}
\newtheorem{definition}{Definition}[section]
\newtheorem{proposition}[definition]{Proposition}
\newtheorem{theorem}[definition]{Theorem}
\newtheorem{corollary}[definition]{Corollary}
\newtheorem{remark}[definition]{Remark}

\newtcolorbox{boardexample}[1]{
  enhanced,breakable,colback=SoftBlue,colframe=CourseBlue,
  boxrule=0.55pt,arc=1mm,left=1.5mm,right=1.5mm,top=1mm,bottom=1mm,
  before skip=0.55em,after skip=0.55em,
  title={Example --- #1},fonttitle=\bfseries}
\newtcolorbox{studentexercise}[1]{
  enhanced,breakable,colback=SoftGold,colframe=CourseGold,
  boxrule=0.55pt,arc=1mm,left=1.5mm,right=1.5mm,top=1mm,bottom=1mm,
  before skip=0.55em,after skip=0.35em,
  title={Exercise --- #1},fonttitle=\bfseries}
\newtcolorbox{shortsolution}[1][]{
  enhanced,breakable,colback=SoftGreen,colframe=CourseGreen,
  boxrule=0.55pt,arc=1mm,left=1.5mm,right=1.5mm,top=1mm,bottom=1mm,
  before skip=0.25em,after skip=0.6em,
  title={Short solution#1},fonttitle=\bfseries}
\newtcolorbox{keyidea}[1]{
  enhanced,breakable,colback=white,colframe=CourseGreen,
  boxrule=0.7pt,arc=1mm,left=1.5mm,right=1.5mm,top=1mm,bottom=1mm,
  before skip=0.5em,after skip=0.5em,
  title={#1},fonttitle=\bfseries}
\newtcolorbox{warningbox}[1]{
  enhanced,breakable,colback=SoftRed,colframe=CourseRed,
  boxrule=0.65pt,arc=1mm,left=1.5mm,right=1.5mm,top=1mm,bottom=1mm,
  before skip=0.5em,after skip=0.5em,
  title={#1},fonttitle=\bfseries}

\title{\textbf{The It\^o Formula: Intuition, Rules, and Applications}\\
\large Lecture 4 Notes --- Class Support and Student Summary}
\author{}
\date{}

\begin{document}
\maketitle
\vspace{-1.4em}

\begin{keyidea}{Goals and logical structure}
Brownian quadratic variation turns the familiar chain rule into a second-order
calculus. The lecture proceeds through one reusable sequence:
\[
\text{quadratic variation}
\longrightarrow
\text{multiplication table}
\longrightarrow
\text{\(\Ito\) formula}
\longrightarrow
\text{SDE transformations and finance}.
\]
By the end, we can transform one- and two-factor \(\Ito\) processes, solve GBM
and OU models, compute the mean of a CIR process, and explain why correlation,
convexity, discounting, and delta hedging all create or remove drift terms.
Exercises are placed immediately after the relevant rule; complete short
solutions follow them.
\end{keyidea}

% ============================================================
\section{Quadratic Variation Becomes a Calculus}
% ============================================================

\subsection{The correction already appeared in Lecture 3}

Lecture 3 established the identity
\[
\int_0^T W_t\dd W_t=\frac12(W_T^2-T),
\]
or equivalently
\[
W_T^2=2\int_0^T W_t\dd W_t+T.
\]
The first term is the one predicted by the ordinary chain rule. The extra
\(T\) is the quadratic variation \([W]_T\).

On one interval,
\[
(W_t+\Delta W)^2-W_t^2
=2W_t\Delta W+(\Delta W)^2.
\]
Brownian scaling gives \(\Delta W=O_{\Pp}(\sqrt{\Delta t})\), so
\((\Delta W)^2=O_{\Pp}(\Delta t)\). The second-order terms are individually
small, but there are order \(1/\Delta t\) of them. Their accumulated effect is
therefore order one.

\begin{center}
\begin{tabular}{lccc}
\toprule
One-step term & Typical size & Sum over \(n\) steps & Limit scale\\
\midrule
\(\Delta t\) & \(n^{-1}\) & \(1\) & survives\\
\((\Delta W)^2\) & \(n^{-1}\) & \(1\) & survives\\
\(\Delta t\,\Delta W\) & \(n^{-3/2}\) & \(n^{-1/2}\) & vanishes\\
\((\Delta t)^2\) & \(n^{-2}\) & \(n^{-1}\) & vanishes\\
\bottomrule
\end{tabular}
\end{center}

\begin{keyidea}{Bookkeeping principle}
Retain every term whose sum can accumulate at order \(\dd t\). For a continuous
diffusion this means keeping first-order terms and the quadratic products of
Brownian increments.
\end{keyidea}

\subsection{Covariation and the differential tables}

\begin{definition}[Quadratic covariation]
For continuous processes \(X,Y\) and a partition
\(\pi=\{0=t_0<\cdots<t_n=t\}\), define
\[
[X,Y]_t^\pi
=\sum_{i=0}^{n-1}
\bigl(X_{t_{i+1}}-X_{t_i}\bigr)
\bigl(Y_{t_{i+1}}-Y_{t_i}\bigr).
\]
When these sums converge in probability as the mesh tends to zero, their
limit is denoted by \([X,Y]_t\). We write \([X]_t=[X,X]_t\).
\end{definition}

The basic examples are
\[
[W]_t=t,\qquad [W,t]_t=0,\qquad [t]_t=0.
\]
The differential notation
\(\dd X_t\dd Y_t:=\dd[X,Y]_t\) summarizes these limiting identities.
For one Brownian driver,
\[
\begin{array}{c|cc}
\toprule
\times & \dd t & \dd W_t\\
\midrule
\dd t & 0 & 0\\
\dd W_t & 0 & \dd t\\
\bottomrule
\end{array}
\qquad\Longleftrightarrow\qquad
(\dd t)^2=0,\quad
\dd t\,\dd W_t=0,\quad
(\dd W_t)^2=\dd t.
\]

\begin{warningbox}{The table is a mnemonic for limits}
These are not literal algebraic infinitesimals. For example,
\((\dd W_t)^2=\dd t\) means that accumulated squared Brownian increments
converge to elapsed time.
\end{warningbox}

Let \(W,B\) be Brownian motions with constant correlation \(\rho\):
\[
[W,B]_t=\rho t,\qquad -1\leq\rho\leq1.
\]
The two-factor table is
\[
\begin{array}{c|ccc}
\toprule
\times & \dd t & \dd W_t & \dd B_t\\
\midrule
\dd t & 0 & 0 & 0\\
\dd W_t & 0 & \dd t & \rho\,\dd t\\
\dd B_t & 0 & \rho\,\dd t & \dd t\\
\bottomrule
\end{array}
\]
and, for a Brownian vector,
\[
\dd W_t^i\dd W_t^j=\rho_{ij}\dd t.
\]

\begin{boardexample}{correlation becomes drift in a product}
The product identity gives
\[
\dd(W_tB_t)
=W_t\dd B_t+B_t\dd W_t+\dd W_t\dd B_t
=W_t\dd B_t+B_t\dd W_t+\rho\dd t.
\]
On a bounded horizon the stochastic integrals have mean zero, so
\[
\E[W_tB_t]=\rho t,
\]
which agrees with the covariance of the jointly Gaussian pair.
\end{boardexample}

\begin{studentexercise}{a first table drill}
Let
\[
\dd X_t=2\dd t+3\dd W_t,\qquad
\dd Y_t=-\dd t+4\dd W_t,\qquad
\dd Z_t=-\dd t+4\dd B_t.
\]
Compute \((\dd X_t)^2\), \(\dd X_t\dd Y_t\), and
\(\dd X_t\dd Z_t\). Treat first the independent case and then the case
\(\dd W_t\dd B_t=-\tfrac12\dd t\).
\end{studentexercise}
\begin{shortsolution}
Every product containing \(\dd t\) vanishes. Therefore
\[
(\dd X_t)^2=9\dd t,\qquad
\dd X_t\dd Y_t=12\dd t,\qquad
\dd X_t\dd Z_t=12\rho\dd t.
\]
The last answer is \(0\) under independence and \(-6\dd t\) when
\(\rho=-1/2\). Only diffusion loadings and correlation enter covariation.
\end{shortsolution}

\begin{studentexercise}{general covariation formula}
Suppose
\[
\dd X_t=a_t\dd t+b_t\dd W_t+c_t\dd B_t,\qquad
\dd Y_t=p_t\dd t+q_t\dd W_t+r_t\dd B_t,
\]
with \(\dd W_t\dd B_t=\rho\dd t\). Find
\(\dd[X,Y]_t\).
\end{studentexercise}
\begin{shortsolution}
\[
\boxed{
\dd[X,Y]_t
=\bigl(b_tq_t+c_tr_t+\rho b_tr_t+\rho c_tq_t\bigr)\dd t.}
\]
The coefficients \(a_t,p_t\) do not appear because finite-variation terms
have zero covariation with continuous semimartingales.
\end{shortsolution}

% ============================================================
\section{The It\^o Formula Is a Second-Order Chain Rule}
% ============================================================

\subsection{Heuristic and precise statement}

For a smooth \(f(t,x)\), Taylor expansion suggests
\[
\Delta f
\approx f_t\Delta t+f_x\Delta W+\frac12f_{xx}(\Delta W)^2.
\]
The mixed term \(\Delta t\Delta W\) and the cubic term
\((\Delta W)^3\) have accumulated order tending to zero. Replacing the
quadratic term by \(\Delta t\) leads to the correct formula.

\begin{theorem}[Brownian \(\Ito\) formula]
Let \(f\in C^{1,2}([0,T]\times\R)\). Then
\[
\begin{aligned}
f(t,W_t)=f(0,W_0)
&+\int_0^t f_t(s,W_s)\dd s
+\int_0^t f_x(s,W_s)\dd W_s\\
&+\frac12\int_0^t f_{xx}(s,W_s)\dd s.
\end{aligned}
\]
Equivalently,
\[
\boxed{
\dd f(t,W_t)
=\left(f_t+\frac12f_{xx}\right)(t,W_t)\dd t
+f_x(t,W_t)\dd W_t.}
\]
\end{theorem}

\begin{proof}[Construction idea]
On a partition \(0=t_0<\cdots<t_n=t\), telescope:
\[
f(t,W_t)-f(0,W_0)=\sum_i
\bigl(f(t_{i+1},W_{t_{i+1}})-f(t_i,W_{t_i})\bigr).
\]
A Taylor expansion at \((t_i,W_{t_i})\) gives
\[
\Delta f_i
=f_t\Delta t_i+f_x\Delta W_i
+\frac12f_{xx}(\Delta W_i)^2+r_i.
\]
The first sum converges to \(\int f_t\dd t\). The left-point second sum
converges to the \(\Ito\) integral \(\int f_x\dd W\). Quadratic variation
turns the third sum into \(\frac12\int f_{xx}\dd t\). After localization,
smoothness controls the derivatives on compact sets and the accumulated
remainder converges to zero in probability.
\end{proof}

\begin{keyidea}{Four-step calculation recipe}
\begin{enumerate}
  \item Treat \(t\) and \(x\) as deterministic variables.
  \item Compute \(f_t(t,x)\), \(f_x(t,x)\), and \(f_{xx}(t,x)\).
  \item Substitute \(x=W_t\).
  \item Group the \(\dd t\) and \(\dd W_t\) terms separately.
\end{enumerate}
Differentiate with respect to the placeholder \(x\), then substitute the
process. This simple notation discipline prevents many errors.
\end{keyidea}

\subsection{First examples and Brownian moments}

\begin{boardexample}{three immediate transforms}
For \(f(x)=x^2\),
\[
\dd(W_t^2)=2W_t\dd W_t+\dd t.
\]
Because
\[
\E\int_0^T(2W_s)^2\dd s
=4\int_0^T s\dd s=2T^2<\infty,
\]
the process \(W_t^2-t\) is a true martingale.

For \(f(x)=x^3\),
\[
\dd(W_t^3)=3W_t^2\dd W_t+3W_t\dd t.
\]
For \(f(x)=\sin x\),
\[
\dd(\sin W_t)
=\cos(W_t)\dd W_t-\frac12\sin(W_t)\dd t.
\]
Even though \(W\) has zero drift, curvature can create drift in \(f(W)\).
\end{boardexample}

\begin{studentexercise}{manufacture a fourth-order martingale}
Apply the formula to \(W_t^4\). Find \(c\) such that
\[
W_t^4-c\int_0^tW_s^2\dd s
\]
has zero drift, and recover \(\E[W_t^4]\).
\end{studentexercise}
\begin{shortsolution}
\[
\dd(W_t^4)=4W_t^3\dd W_t+6W_t^2\dd t.
\]
Thus
\[
M_t=W_t^4-6\int_0^tW_s^2\dd s
\]
has zero drift. Since
\(\E\int_0^T16W_s^6\dd s<\infty\), it is a true martingale on
\([0,T]\). Hence
\[
\E[W_t^4]
=6\int_0^t\E[W_s^2]\dd s
=6\int_0^t s\dd s
=3t^2.
\]
\end{shortsolution}

\begin{studentexercise}{one more step gives the sixth moment}
Use \(\E[W_t^4]=3t^2\) and the \(\Ito\) formula for \(x^6\) to compute
\(\E[W_T^6]\).
\end{studentexercise}
\begin{shortsolution}
\[
\dd(W_t^6)=6W_t^5\dd W_t+15W_t^4\dd t.
\]
The integrand is square-integrable on bounded horizons because Brownian
motion has finite moments of every order. Taking expectations,
\[
\boxed{
\E[W_T^6]
=15\int_0^T\E[W_t^4]\dd t
=45\int_0^Tt^2\dd t
=15T^3.}
\]
More generally, \(m_n(t)=\E[W_t^n]\) satisfies
\[
m_n'(t)=\frac{n(n-1)}2m_{n-2}(t).
\]
\end{shortsolution}

\subsection{Stopping turns martingales into exit-time identities}

For \(a>0\), define the first exit time
\[
\tau_a=\inf\{t\geq0:|W_t|=a\}.
\]
The safe order of operations is to stop first at the bounded time
\(\tau_t=t\wedge\tau_a\), apply optional sampling, and only then let
\(t\to\infty\).

\begin{studentexercise}{exit time and occupation inside an interval}
Use
\[
W_t^2-t
\qquad\text{and}\qquad
W_t^4-6\int_0^tW_s^2\dd s
\]
to determine
\[
\E[\tau_a],
\qquad
\E\left[\int_0^{\tau_a}W_s^2\dd s\right].
\]
\end{studentexercise}
\begin{shortsolution}
Optional sampling at \(\tau_t\) gives
\[
\E[W_{\tau_t}^2]=\E[\tau_t]\leq a^2.
\]
Monotone convergence yields \(\E[\tau_a]\leq a^2\), hence
\(\tau_a<\infty\) almost surely. Brownian continuity gives
\(|W_{\tau_a}|=a\). Since \(|W_{\tau_t}|\leq a\), dominated convergence
then yields
\[
\boxed{\E[\tau_a]=a^2.}
\]
Stopping the fourth-order martingale gives
\[
\E[W_{\tau_t}^4]
=6\E\left[\int_0^{\tau_t}W_s^2\dd s\right].
\]
Dominated convergence on the left and monotone convergence on the right give
\[
\boxed{
\E\left[\int_0^{\tau_a}W_s^2\dd s\right]=\frac{a^4}{6}.}
\]
\end{shortsolution}

\subsection{Exponentials and the backward heat equation}

\begin{proposition}[Exponential martingale]
For \(\lambda\in\R\),
\[
M_t=\exp\left(\lambda W_t-\frac12\lambda^2t\right)
\]
satisfies
\[
\dd M_t=\lambda M_t\dd W_t.
\]
Moreover \(M_t\geq0\) and a direct Gaussian calculation gives
\(\E[M_t]=1\), so \(M\) is a true martingale.
\end{proposition}

\begin{proof}
For \(f(t,x)=e^{\lambda x-\lambda^2t/2}\),
\[
f_t=-\frac12\lambda^2f,\qquad
f_x=\lambda f,\qquad
f_{xx}=\lambda^2f.
\]
The two drift terms cancel in the \(\Ito\) formula. Also,
\[
\E[M_t]
=e^{-\lambda^2t/2}\E[e^{\lambda W_t}]
=e^{-\lambda^2t/2}e^{\lambda^2t/2}=1.
\]
\end{proof}

More generally, for \(Y_t=u(t,W_t)\),
\[
\dd Y_t
=\left(u_t+\frac12u_{xx}\right)(t,W_t)\dd t
+u_x(t,W_t)\dd W_t.
\]
Thus the drift vanishes when
\[
\boxed{u_t+\frac12u_{xx}=0.}
\]
This is the backward heat equation. Suitable growth or integrability
conditions are still needed to upgrade the resulting local martingale to a
true martingale.

\begin{studentexercise}{which transforms have zero drift?}
For fixed \(T>0\), test
\[
u_1(t,x)=x^2-t,\quad
u_2(t,x)=x^2+t,\quad
u_3(t,x)=x^2+T-t,\quad
u_4(t,x)=e^{\lambda x-\lambda^2t/2}.
\]
Compute \(u_t+\tfrac12u_{xx}\) and interpret \(u_3(t,W_t)\).
\end{studentexercise}
\begin{shortsolution}
\[
\begin{array}{lcc}
\toprule
u & u_t+\frac12u_{xx} & \text{conclusion}\\
\midrule
x^2-t & 0 & \text{martingale}\\
x^2+t & 2 & \text{drift }2\dd t\\
x^2+T-t & 0 & \text{martingale}\\
e^{\lambda x-\lambda^2t/2} & 0 & \text{martingale}\\
\bottomrule
\end{array}
\]
Since the third process equals \(W_T^2\) at \(t=T\),
\[
W_t^2+T-t=\E[W_T^2\mid\cF_t].
\]
\end{shortsolution}

\begin{studentexercise}{a conditional fourth moment}
Find a polynomial \(u(t,x)\) such that
\[
u(T,x)=x^4,\qquad u_t+\frac12u_{xx}=0.
\]
Deduce a formula for \(\E[W_T^4\mid\cF_t]\).
\end{studentexercise}
\begin{shortsolution}
Let \(h=T-t\). Direct differentiation shows that
\[
u(t,x)=x^4+6hx^2+3h^2
\]
solves the backward heat equation and has terminal value \(x^4\). Therefore
\[
\boxed{
\E[W_T^4\mid\cF_t]
=W_t^4+6(T-t)W_t^2+3(T-t)^2.}
\]
This also follows by expanding
\(W_T=W_t+(W_T-W_t)\) and using Gaussian moments.
\end{shortsolution}

% ============================================================
\section{From Brownian Motion to It\^o Processes}
% ============================================================

\subsection{Definition, quadratic variation, and the general formula}

\begin{definition}[\(\Ito\) process]
On a fixed horizon \([0,T]\), an \(\Ito\) process has the representation
\[
X_t=X_0+\int_0^t a_s\dd s+\int_0^t b_s\dd W_s,
\]
where, for example, \(a,b\) are progressively measurable and
\[
\int_0^T|a_s|\dd s<\infty\quad\text{a.s.},
\qquad
\E\left[\int_0^Tb_s^2\dd s\right]<\infty.
\]
We use the shorthand
\[
\dd X_t=a_t\dd t+b_t\dd W_t.
\]
\end{definition}

\begin{keyidea}{Differential form versus integral form}
The differential notation is a compact way to record the local drift and
diffusion coefficients. Its precise meaning is the integral identity
\[
\boxed{
\dd X_t=a_t\dd t+b_t\dd W_t
\quad\Longleftrightarrow\quad
X_t-X_s=\int_s^t a_u\dd u+\int_s^t b_u\dd W_u,}
\]
valid for every \(0\leq s\leq t\leq T\) (up to the usual almost-sure
qualification). For example,
\[
\dd X_t=(m-X_t)\dd t+\sigma\dd W_t
\]
means
\[
X_t=X_0+\int_0^t(m-X_u)\dd u+\sigma W_t.
\]
The differential form is convenient for It\^o calculations; the integral
form is the actual equality defining the process. In particular, one should
not divide by \(\dd t\): Brownian paths are not differentiable, so
\(\dd W_t/\dd t\) is not an ordinary random function.
\end{keyidea}

\begin{proposition}[Uniqueness of the drift--noise decomposition]
Fix the filtration and the Brownian motion \(W\). Suppose that the same
continuous process admits two representations
\[
\begin{aligned}
X_t
&=X_0+\int_0^t a_s\dd s+\int_0^t b_s\dd W_s\\
&=X_0+\int_0^t\widetilde a_s\dd s
       +\int_0^t\widetilde b_s\dd W_s,
\end{aligned}
\]
under the integrability assumptions in the preceding definition. Then
\[
\boxed{
a=\widetilde a,\qquad b=\widetilde b,
\qquad \dd t\otimes\dd\Pp\text{-almost everywhere}.}
\]
\end{proposition}

\begin{proof}
Subtract the two decompositions and set
\[
A_t:=\int_0^t(a_s-\widetilde a_s)\dd s,
\qquad
M_t:=\int_0^t(b_s-\widetilde b_s)\dd W_s.
\]
Then \(A_t+M_t=0\) for every \(t\). The process \(A\) has continuous
finite-variation paths, whereas \(M\) is a continuous local martingale.
Consequently, \(M=-A\) is a continuous local martingale of finite variation.
Such a process is constant. Since \(A_0=M_0=0\), we obtain
\[
A_t=M_t=0
\qquad\text{for every }t,\quad\text{a.s.}
\]
The quadratic variation of \(M\) is therefore zero:
\[
0=[M]_t
=\int_0^t(b_s-\widetilde b_s)^2\dd s.
\]
It follows that \(b=\widetilde b\),
\(\dd t\otimes\dd\Pp\)-almost everywhere. Finally, \(A\equiv0\) and the
absolute continuity of \(A\) imply
\(a=\widetilde a\), \(\dd t\otimes\dd\Pp\)-almost everywhere.
\end{proof}

\begin{remark}
The underlying principle is the uniqueness of the canonical continuous
semimartingale decomposition into a finite-variation part and a local
martingale part. Equality of the coefficients is an almost-everywhere
statement; changing \(a\) or \(b\) on a
\(\dd t\otimes\dd\Pp\)-null set does not change \(X\).
\end{remark}

The drift \(a_t\dd t\) is the local trend; the diffusion
\(b_t\dd W_t\) is the local random shock. Applying the multiplication table,
\[
(\dd X_t)^2
=a_t^2(\dd t)^2+2a_tb_t\dd t\dd W_t+b_t^2(\dd W_t)^2
=b_t^2\dd t.
\]
Hence
\[
\boxed{
\dd[X]_t=b_t^2\dd t,\qquad
[X]_t=\int_0^tb_s^2\dd s.}
\]
Quadratic variation records volatility, not drift.

\begin{theorem}[One-dimensional \(\Ito\) formula]
Let
\[
\dd X_t=a_t\dd t+b_t\dd W_t,
\qquad f\in C^{1,2}([0,T]\times\R).
\]
Then
\[
\boxed{
\dd f(t,X_t)
=\left(f_t+a_tf_x+\frac12b_t^2f_{xx}\right)(t,X_t)\dd t
+b_tf_x(t,X_t)\dd W_t.}
\]
Equivalently,
\[
\dd f=f_t\dd t+f_x\dd X+\frac12f_{xx}\dd[X].
\]
\end{theorem}

\begin{boardexample}{the square of a general diffusion}
For \(f(x)=x^2\),
\[
\boxed{
\dd(X_t^2)
=(2X_ta_t+b_t^2)\dd t+2X_tb_t\dd W_t.}
\]
Whenever the stochastic integral is a true martingale,
\[
X_t^2-X_0^2-\int_0^t(2X_sa_s+b_s^2)\dd s
\]
is a martingale.
\end{boardexample}

\begin{studentexercise}{arithmetic diffusion and its second moment}
Let
\[
X_t=x+\mu t+\sigma W_t.
\]
Compute \(\dd(X_t^2)\), find a compensated process with zero drift, and
calculate \(\E[X_t^2]\).
\end{studentexercise}
\begin{shortsolution}
\[
\dd(X_t^2)
=(2\mu X_t+\sigma^2)\dd t+2\sigma X_t\dd W_t.
\]
The stochastic integrand is square-integrable on bounded horizons because
\(\int_0^T\E[X_s^2]\dd s<\infty\). Taking expectations and either solving
the resulting integral equation or using the known Gaussian law gives
\[
\boxed{\E[X_t^2]=(x+\mu t)^2+\sigma^2t.}
\]
This is \((\text{mean})^2+\text{variance}\).
\end{shortsolution}

\begin{warningbox}{When may a stochastic integral be erased in expectation?}
If
\[
\E\left[\int_0^TH_s^2\dd s\right]<\infty,
\]
then \(\int_0^tH_s\dd W_s\) is an \(L^2\)-martingale and has expectation
zero. Pathwise local square-integrability normally gives only a local
martingale. Zero drift alone is not enough. Absolute integrability is also
needed before exchanging expectation with a drift integral.
\end{warningbox}

\subsection{Logarithms and inverses}

\begin{proposition}[Reusable log rule]
Suppose \(X_t>0\) and
\[
\frac{\dd X_t}{X_t}=\alpha_t\dd t+\beta_t\dd W_t.
\]
Then
\[
\boxed{
\dd\log X_t
=\left(\alpha_t-\frac12\beta_t^2\right)\dd t
+\beta_t\dd W_t.}
\]
\end{proposition}

\begin{proof}
For \(f(x)=\log x\),
\[
f'(x)=\frac1x,\qquad f''(x)=-\frac1{x^2}.
\]
Substitution into the general formula gives the claim. The term
\(-\beta_t^2/2\) is often called volatility drag.
\end{proof}

If strict positivity is not known globally, the identity is valid only up to
the first hitting time of zero.

\begin{studentexercise}{invert an \(\Ito\) process}
Let \(X_t\neq0\) and
\[
\dd X_t=a_t\dd t+b_t\dd W_t.
\]
Find \(\dd(1/X_t)\), then specialize to
\(\dd X_t/X_t=\mu\dd t+\sigma\dd W_t\).
\end{studentexercise}
\begin{shortsolution}
With \(f(x)=x^{-1}\),
\[
f'(x)=-x^{-2},\qquad f''(x)=2x^{-3}.
\]
Therefore
\[
\boxed{
\dd\left(\frac1{X_t}\right)
=\left(-\frac{a_t}{X_t^2}+\frac{b_t^2}{X_t^3}\right)\dd t
-\frac{b_t}{X_t^2}\dd W_t.}
\]
For geometric dynamics and \(R_t=1/X_t\),
\[
\boxed{
\frac{\dd R_t}{R_t}
=(\sigma^2-\mu)\dd t-\sigma\dd W_t.}
\]
The \(+\sigma^2\) term is the correction that an ordinary quotient rule
would miss. The product rule applied to \(X_tR_t=1\) verifies the answer.
\end{shortsolution}

% ============================================================
\section{Transformations Solve Canonical SDEs}
% ============================================================

\subsection{Geometric Brownian motion}

Consider
\[
\dd S_t=\mu S_t\dd t+\sigma S_t\dd W_t,
\qquad S_0>0,\quad \sigma>0.
\]
The logarithm converts relative shocks into additive shocks:
\[
\dd\log S_t
=\left(\mu-\frac12\sigma^2\right)\dd t+\sigma\dd W_t.
\]
Integration yields the explicit solution
\[
\boxed{
S_t=S_0\exp\left\{
\left(\mu-\frac12\sigma^2\right)t+\sigma W_t
\right\}.}
\]
It is strictly positive, which justifies the logarithmic transformation.
Moreover,
\[
\log S_t\sim
N\left(
\log S_0+\left(\mu-\frac12\sigma^2\right)t,\,
\sigma^2t
\right),
\]
so \(S_t\) is lognormal.

\begin{proposition}[Mean, median, and a target probability]
\[
\E[S_t]=S_0e^{\mu t},
\qquad
\Med(S_t)=S_0e^{(\mu-\sigma^2/2)t}.
\]
For \(K>0\),
\[
\Pp(S_T>K)
=1-\Phi\left(
\frac{\log(K/S_0)-(\mu-\sigma^2/2)T}
{\sigma\sqrt T}
\right).
\]
\end{proposition}

\begin{boardexample}{one year of GBM in numbers}
Let
\[
S_0=100,\qquad \mu=8\%,\qquad \sigma=20\%,\qquad T=1.
\]
Then
\[
\E[S_1]=100e^{0.08}\approx108.33,
\qquad
\Med(S_1)=100e^{0.06}\approx106.18.
\]
For the target \(K=110\),
\[
\Pp(S_1>110)
=1-\Phi\left(\frac{\log(1.10)-0.06}{0.20}\right)
\approx43.0\%.
\]
The mean, median, and target probability answer different questions. A
positive expected return does not eliminate downside uncertainty.
\end{boardexample}

\begin{studentexercise}{powers of a GBM}
For \(p\in\R\), let \(Y_t=S_t^p\).
\begin{enumerate}
  \item Find \(\dd Y_t/Y_t\).
  \item Choose \(k\) such that \(e^{-kt}S_t^p\) has zero drift.
  \item Check \(p=0,1,2,-1\).
\end{enumerate}
\end{studentexercise}
\begin{shortsolution}
\[
\boxed{
\frac{\dd(S_t^p)}{S_t^p}
=\left[p\mu+\frac12p(p-1)\sigma^2\right]\dd t
+p\sigma\dd W_t.}
\]
Thus
\[
k=p\mu+\frac12p(p-1)\sigma^2.
\]
For constant coefficients, \(e^{-kt}S_t^p\) is a constant multiple of an
exponential Brownian martingale and is therefore a true martingale. For
\(p=-1\), the drift is \(\sigma^2-\mu\), matching the inverse calculation.
\end{shortsolution}

If coefficients vary through time,
\[
\frac{\dd S_t}{S_t}
=\mu_t\dd t+\sigma_t\dd W_t,
\]
the same log transform gives
\[
\boxed{
S_t=S_0\exp\left\{
\int_0^t\left(\mu_s-\frac12\sigma_s^2\right)\dd s
+\int_0^t\sigma_s\dd W_s
\right\}.}
\]
This holds under the usual progressive measurability and local integrability
conditions. If \(\mu,\sigma\) are deterministic, the log price remains
Gaussian.

\subsection{Ornstein--Uhlenbeck and Vasicek mean reversion}

The OU SDE is
\[
\dd X_t=\kappa(m-X_t)\dd t+\sigma\dd W_t,
\qquad \kappa>0.
\]
Above \(m\), the drift is negative; below \(m\), it is positive. Multiplying
by \(e^{\kappa t}\) removes the linear drift. Because this integrating factor
has finite variation,
\[
\begin{aligned}
\dd(e^{\kappa t}X_t)
&=e^{\kappa t}\dd X_t+\kappa e^{\kappa t}X_t\dd t\\
&=\kappa m e^{\kappa t}\dd t+\sigma e^{\kappa t}\dd W_t.
\end{aligned}
\]
Integrating between \(s\) and \(t\) gives
\[
\boxed{
X_t=m+(X_s-m)e^{-\kappa(t-s)}
+\sigma\int_s^t e^{-\kappa(t-u)}\dd W_u.}
\]
Conditionally on \(\cF_s\), the stochastic integral is centered Gaussian.
Therefore
\[
\E[X_t\mid\cF_s]
=m+(X_s-m)e^{-\kappa(t-s)}
\]
and
\[
\Var(X_t\mid\cF_s)
=\sigma^2\int_s^te^{-2\kappa(t-u)}\dd u
=\frac{\sigma^2}{2\kappa}
\left(1-e^{-2\kappa(t-s)}\right).
\]

\begin{studentexercise}{read the long-run behavior from the OU formula}
Find the limiting conditional mean, limiting conditional variance, and
half-life of the deviation from \(m\). Evaluate the half-life at
\(\kappa=0.7\) per year.
\end{studentexercise}
\begin{shortsolution}
\[
\E[X_t\mid\cF_s]\longrightarrow m,
\qquad
\Var(X_t\mid\cF_s)\longrightarrow\frac{\sigma^2}{2\kappa}.
\]
The half-life \(h\) solves \(e^{-\kappa h}=1/2\), hence
\[
\boxed{h=\frac{\log2}{\kappa}.}
\]
When \(\kappa=0.7\), \(h\approx0.99\) year. After two years,
\[
e^{-0.7\times2}\approx24.7\%
\]
of the current deviation remains in the conditional mean.
\end{shortsolution}

\subsection{The first moment of a CIR process}

Assume a nonnegative process satisfies
\[
\dd R_t=(\alpha-\beta R_t)\dd t
+\sigma\sqrt{R_t}\dd W_t,
\qquad
R_0=r_0,\quad \alpha,\beta,\sigma>0.
\]
The integrating factor \(e^{\beta t}\) gives
\[
\dd(e^{\beta t}R_t)
=\alpha e^{\beta t}\dd t
+\sigma e^{\beta t}\sqrt{R_t}\dd W_t.
\]
After localization and the usual integrability check, the stochastic
integral has mean zero. Hence
\[
\E[e^{\beta t}R_t]
=r_0+\frac{\alpha}{\beta}(e^{\beta t}-1),
\]
so
\[
\boxed{
\E[R_t]
=r_0e^{-\beta t}
+\frac{\alpha}{\beta}(1-e^{-\beta t})
\longrightarrow\frac{\alpha}{\beta}.}
\]
The square-root volatility controls dispersion, but it does not enter the
first-moment ODE.

\begin{studentexercise}{optional extension: the CIR second moment}
Apply the \(\Ito\) formula to \(R_t^2\). Obtain an ODE for
\(m_2(t)=\E[R_t^2]\) in terms of \(m_1(t)=\E[R_t]\), and find the limiting
variance.
\end{studentexercise}
\begin{shortsolution}
\[
\dd(R_t^2)
=\left((2\alpha+\sigma^2)R_t-2\beta R_t^2\right)\dd t
+2\sigma R_t^{3/2}\dd W_t.
\]
Thus
\[
m_2'(t)+2\beta m_2(t)=(2\alpha+\sigma^2)m_1(t).
\]
Using the integrating factor \(e^{2\beta t}\),
\[
\begin{aligned}
m_2(t)
={}&r_0^2e^{-2\beta t}\\
&+(2\alpha+\sigma^2)
\left[
\frac{\alpha}{2\beta^2}(1-e^{-2\beta t})
+\frac{r_0-\alpha/\beta}{\beta}
(e^{-\beta t}-e^{-2\beta t})
\right].
\end{aligned}
\]
Therefore
\[
\lim_{t\to\infty}\Var(R_t)
=\frac{\alpha\sigma^2}{2\beta^2}.
\]
\end{shortsolution}

% ============================================================
\section{Products, Ratios, and Several Risk Factors}
% ============================================================

\subsection{The multidimensional formula and product rule}

\begin{theorem}[Two-dimensional \(\Ito\) formula]
Let \(f\in C^{1,2}([0,T]\times\R^2)\) and let \(X,Y\) be continuous
\(\Ito\) processes. Then
\[
\begin{aligned}
\dd f(t,X_t,Y_t)
={}&f_t\dd t+f_x\dd X_t+f_y\dd Y_t\\
&+\frac12f_{xx}\dd[X]_t
+f_{xy}\dd[X,Y]_t
+\frac12f_{yy}\dd[Y]_t.
\end{aligned}
\]
\end{theorem}

The mixed derivative has no factor \(1/2\): it appears twice in the symmetric
Hessian sum. If the Brownian drivers have correlation \(\rho\), the mixed
term carries that correlation.

Choosing \(f(x,y)=xy\) gives the stochastic product rule
\[
\boxed{
\dd(X_tY_t)
=X_t\dd Y_t+Y_t\dd X_t+\dd[X,Y]_t.}
\]
For
\[
\dd X_t=a_t\dd t+b_t\dd W_t,\qquad
\dd Y_t=c_t\dd t+d_t\dd B_t,
\]
we have
\[
\dd[X,Y]_t=\rho b_td_t\dd t.
\]

\begin{boardexample}{product of two correlated assets}
Suppose
\[
\frac{\dd S_t}{S_t}
=\mu_S\dd t+\sigma_S\dd W_t,\qquad
\frac{\dd V_t}{V_t}
=\mu_V\dd t+\sigma_V\dd B_t,
\]
and \(\dd W_t\dd B_t=\rho\dd t\). Then
\[
\boxed{
\frac{\dd(S_tV_t)}{S_tV_t}
=\left(\mu_S+\mu_V+\rho\sigma_S\sigma_V\right)\dd t
+\sigma_S\dd W_t+\sigma_V\dd B_t.}
\]
Positive correlation increases the product drift through covariation;
negative correlation decreases it.
\end{boardexample}

\begin{studentexercise}{normalize a sum of Brownian motions}
Let \(W,B\) be jointly Brownian with correlation \(\rho>-1\). Find a
constant \(c\) such that \(c(W+B)\) is a standard Brownian motion. What
happens at \(\rho=-1\)?
\end{studentexercise}
\begin{shortsolution}
\[
\dd[W+B]_t
=(\dd W_t+\dd B_t)^2
=2(1+\rho)\dd t.
\]
Hence
\[
\boxed{
\widetilde W_t=\frac{W_t+B_t}{\sqrt{2(1+\rho)}}}
\]
has quadratic variation \(t\). Joint Gaussianity, or Levy's characterization,
shows that it is Brownian. At \(\rho=-1\), \(B=-W\) up to
indistinguishability and the sum is identically zero.
\end{shortsolution}

\subsection{Ratios and inverse corrections}

\begin{studentexercise}{ratio of two risky assets}
With the asset dynamics above, let \(R_t=S_t/V_t\). Find
\(\dd R_t/R_t\) using \(R=S(1/V)\).
\end{studentexercise}
\begin{shortsolution}
First apply the inverse rule:
\[
\frac{\dd(1/V_t)}{1/V_t}
=(\sigma_V^2-\mu_V)\dd t-\sigma_V\dd B_t.
\]
Then use the product rule:
\[
\begin{aligned}
\frac{\dd R_t}{R_t}
={}&\frac{\dd S_t}{S_t}
+\frac{\dd(1/V_t)}{1/V_t}
+\frac{\dd S_t}{S_t}
  \frac{\dd(1/V_t)}{1/V_t}\\
={}&\left(
\mu_S-\mu_V+\sigma_V^2-\rho\sigma_S\sigma_V
\right)\dd t\\
&+\sigma_S\dd W_t-\sigma_V\dd B_t.
\end{aligned}
\]
Thus
\[
\boxed{
\frac{\dd R_t}{R_t}
=\left(
\mu_S-\mu_V+\sigma_V^2-\rho\sigma_S\sigma_V
\right)\dd t
+\sigma_S\dd W_t-\sigma_V\dd B_t.}
\]
If \(S=V\), every term vanishes. If \(V\) is deterministic,
\(\sigma_V=0\), ordinary quotient calculus returns.
\end{shortsolution}

\begin{studentexercise}{log ratio as a consistency check}
Compute \(\dd\log(S_t/V_t)\) directly from the two log-price dynamics.
Why does \(\rho\) not appear in its drift, even though it appears in the
level-ratio drift?
\end{studentexercise}
\begin{shortsolution}
\[
\boxed{
\dd\log\frac{S_t}{V_t}
=\left[
\mu_S-\mu_V-\frac12(\sigma_S^2-\sigma_V^2)
\right]\dd t
+\sigma_S\dd W_t-\sigma_V\dd B_t.}
\]
Correlation affects the quadratic variation of the final diffusion:
\[
\dd\left[\log(S/V)\right]_t
=\left(\sigma_S^2+\sigma_V^2
-2\rho\sigma_S\sigma_V\right)\dd t.
\]
Applying the exponential rule to the log ratio recreates the
\(-\rho\sigma_S\sigma_V\) term in the level-ratio drift.
\end{shortsolution}

\subsection{Minimum-variance hedging}

Consider the instantaneous shock
\[
\sigma_X\dd W_t-h\sigma_Y\dd B_t,
\qquad \sigma_Y>0,
\qquad
\dd W_t\dd B_t=\rho\dd t.
\]
Its variance rate is
\[
q(h)=\sigma_X^2+h^2\sigma_Y^2
-2h\rho\sigma_X\sigma_Y.
\]
Differentiating,
\[
q'(h)=2h\sigma_Y^2-2\rho\sigma_X\sigma_Y.
\]
Since \(q''(h)=2\sigma_Y^2>0\), the unique minimizer is
\[
\boxed{
h^*=\rho\frac{\sigma_X}{\sigma_Y},
\qquad
q(h^*)=\sigma_X^2(1-\rho^2).}
\]
Only the correlated component can be hedged; orthogonal noise remains.

\begin{boardexample}{numerical hedge ratio}
If
\[
\sigma_X=25\%,\qquad \sigma_Y=20\%,\qquad \rho=0.8,
\]
then
\[
h^*=0.8\frac{0.25}{0.20}=1,
\qquad
\sqrt{q(h^*)}=25\%\sqrt{1-0.8^2}=15\%.
\]
The hedge removes the correlated component in this simplified instantaneous
calculation, but it cannot remove the \(15\%\) orthogonal component.
\end{boardexample}

\subsection{Optional extension: predictable rotations}

Let \(W\) be a standard \(d\)-dimensional Brownian motion and \(H_t\) a
predictable matrix satisfying
\[
H_tH_t^\top=I_d.
\]
Define
\[
B_t=\int_0^tH_s\dd W_s.
\]
Then \(B\) is a continuous local martingale and
\[
\dd[B^i,B^j]_t
=(H_tH_t^\top)_{ij}\dd t
=\delta_{ij}\dd t.
\]
Levy's characterization implies that \(B\) is again a standard
\(d\)-dimensional Brownian motion. In one dimension,
\[
\int_0^t\operatorname{sgn}(W_s)\dd W_s
\]
is Brownian even though its integrand depends on the observed path.

\section{Finance Applications of It\^o Calculus}
% ============================================================

\subsection{Discounting and the probability measure}

Let the bank account and risky asset satisfy
\[
\dd A_t=rA_t\dd t,\qquad
\dd S_t=\mu S_t\dd t+\sigma S_t\dd W_t^{\Pp}.
\]
Since \(A_t^{-1}=A_0^{-1}e^{-rt}\) has finite variation,
\[
\dd(A_t^{-1})=-rA_t^{-1}\dd t
\]
and \(\dd[A^{-1},S]_t=0\). For
\(\widetilde S_t=S_t/A_t\),
\[
\boxed{
\frac{\dd\widetilde S_t}{\widetilde S_t}
=(\mu-r)\dd t+\sigma\dd W_t^{\Pp}.}
\]
Discounting by the bank account does not generally remove the physical drift.

Under a pricing measure \(\Q\), the model is written
\[
\frac{\dd S_t}{S_t}
=r\dd t+\sigma\dd W_t^{\Q},
\]
so
\[
\frac{\dd\widetilde S_t}{\widetilde S_t}
=\sigma\dd W_t^{\Q}.
\]
The discounted price is then a local martingale; standard integrability
conditions make it a true martingale.

\begin{warningbox}{Always state the measure}
The same asset can have drift \(\mu\) under the physical measure and drift
\(r\) under a pricing measure. The martingale property of discounted prices
is a risk-neutral statement, not an automatic physical-measure identity.
\end{warningbox}

\subsection{Option dynamics and delta hedging}

Suppose \(V_t=v(t,S_t)\), with
\[
\dd S_t=\mu S_t\dd t+\sigma S_t\dd W_t.
\]
The \(\Ito\) formula gives
\[
\boxed{
\dd V_t
=\left(
v_t+\mu S_tv_s+\frac12\sigma^2S_t^2v_{ss}
\right)\dd t
+\sigma S_tv_s\dd W_t.}
\]
The first two state derivatives have financial names:
\[
\Delta=v_s,\qquad \Gamma=v_{ss}.
\]
Delta is the first-order price exposure; gamma is curvature.

Choose \(\Delta_t=v_s(t,S_t)\). If changes in the stock holding are financed
through the cash account, the self-financing hedged gain is
\[
\dd\Pi_t:=\dd V_t-\Delta_t\dd S_t.
\]
Substitution gives
\[
\boxed{
\dd\Pi_t
=\left(v_t+\frac12\sigma^2S_t^2v_{ss}\right)\dd t.}
\]
Both the instantaneous Brownian shock and the unknown physical drift
\(\mu\) cancel.

\begin{warningbox}{A bookkeeping point about dynamic delta}
Because \(\Delta_t\) changes through time, it is generally false that
\[
\dd(V_t-\Delta_tS_t)=\dd V_t-\Delta_t\dd S_t.
\]
The right-hand side is the gain of a dynamically financed strategy. Changes
in the cost of the stock position are offset by the cash account.
\end{warningbox}

\begin{boardexample}{curvature creates hedged profit and loss}
For the toy claim \(V_t=S_t^2\),
\[
\Delta_t=2S_t,\qquad \Gamma_t=2.
\]
The \(\Ito\) formula gives
\[
\dd(S_t^2)=2S_t\dd S_t+\sigma^2S_t^2\dd t.
\]
Therefore
\[
\boxed{
\dd V_t-2S_t\dd S_t=\sigma^2S_t^2\dd t.}
\]
Continuous rebalancing converts quadratic price variation into a
finite-variation term.
\end{boardexample}

\begin{boardexample}{one-day gamma scale}
For \(\Gamma=0.02\), \(S=100\), \(\sigma=20\%\), and
\(\Delta t=1/252\),
\[
\frac12\Gamma\sigma^2S^2\Delta t
=\frac12(0.02)(0.20)^2(100)^2\frac1{252}
\approx \$0.016
\]
per option, before theta and financing. The number is small for one day but
accumulates and is highly sensitive to gamma, volatility, and the asset level.
\end{boardexample}

\subsection{Preview of the Black--Scholes equation}

If the delta-hedged position is locally riskless, no-arbitrage requires it to
earn the bank-account rate on its value \(V_t-\Delta_tS_t\):
\[
\dd V_t-\Delta_t\dd S_t
=r(V_t-\Delta_tS_t)\dd t.
\]
Using \(\Delta=v_s\) gives
\[
\boxed{
v_t+\frac12\sigma^2s^2v_{ss}+rsv_s-rv=0.}
\]
Quadratic variation supplies the gamma term; no-arbitrage determines the
growth rate of the locally riskless position. The full pricing argument is
the subject of Lecture 5.

\subsection{Convexity and realized variance}

For \(f\) with no explicit time dependence and
\[
\dd X_t=a_t\dd t+b_t\dd W_t,
\]
the transformed drift is
\[
f'(X_t)a_t+\frac12b_t^2f''(X_t).
\]
Hence
\[
\begin{array}{lcc}
\toprule
\text{shape of }f & \text{sign of }f'' & \text{volatility correction}\\
\midrule
\text{convex} & f''\geq0 & \text{nonnegative}\\
\text{linear} & f''=0 & 0\\
\text{concave} & f''\leq0 & \text{nonpositive}\\
\bottomrule
\end{array}
\]
Volatility raises the drift of a convex transform and lowers the drift of a
concave transform. The examples \(x^2\) and \(\log x\) display the two signs.

For a positive asset,
\[
\frac{\dd S_t}{S_t}
=\mu_t\dd t+\sigma_t\dd W_t,
\qquad
\dd\log S_t
=\left(\mu_t-\frac12\sigma_t^2\right)\dd t
+\sigma_t\dd W_t.
\]
Consequently,
\[
\boxed{
[\log S]_T=\int_0^T\sigma_t^2\dd t.}
\]
Squared high-frequency log returns estimate this integrated variance.

\begin{boardexample}{four-day realized variance}
For log returns
\[
1.0\%,\quad -0.5\%,\quad 0.7\%,\quad -0.2\%,
\]
the realized variance is
\[
\widehat{RV}
=0.01^2+(-0.005)^2+0.007^2+(-0.002)^2
=1.78\times10^{-4}.
\]
Annualizing over four trading days,
\[
\boxed{
\widehat\sigma_{\mathrm{ann}}
=\sqrt{\frac{252}{4}\widehat{RV}}
\approx10.6\%.}
\]
\end{boardexample}

% ============================================================
\section{Guided Practice and Common Traps}
% ============================================================

\subsection{Rapid practice}

\begin{studentexercise}{three transformations}
Compute:
\begin{enumerate}
  \item \(\dd Y_t\) for \(Y_t=e^{\alpha t+\beta W_t}\);
  \item \(\dd Z_t\) for \(Z_t=W_tB_t\) and
        \(\dd W_t\dd B_t=\rho\dd t\);
  \item an expression without a stochastic integral for
        \(\int_0^te^{W_s}\dd W_s\).
\end{enumerate}
\end{studentexercise}
\begin{shortsolution}
\[
\dd Y_t
=Y_t\left(\alpha+\frac12\beta^2\right)\dd t
+\beta Y_t\dd W_t,
\]
\[
\dd Z_t
=W_t\dd B_t+B_t\dd W_t+\rho\dd t,
\]
and
\[
\boxed{
\int_0^te^{W_s}\dd W_s
=e^{W_t}-1-\frac12\int_0^te^{W_s}\dd s.}
\]
The last identity comes from integrating
\(\dd(e^{W_t})=e^{W_t}\dd W_t+\tfrac12e^{W_t}\dd t\).
\end{shortsolution}

\begin{studentexercise}{manufacture a trigonometric martingale}
Let
\[
M_t=e^{at}\cos(bW_t).
\]
Choose \(a\) to remove the drift, then verify that \(M\) is a true
martingale on every bounded horizon.
\end{studentexercise}
\begin{shortsolution}
\[
\dd M_t
=e^{at}\left(a-\frac12b^2\right)\cos(bW_t)\dd t
-be^{at}\sin(bW_t)\dd W_t.
\]
Thus \(a=b^2/2\). On \([0,T]\),
\[
\E\int_0^Tb^2e^{b^2t}\sin^2(bW_t)\dd t
\leq b^2\int_0^Te^{b^2t}\dd t<\infty.
\]
The stochastic integral is therefore in \(L^2\), so \(M\) is a true
martingale.
\end{shortsolution}

\begin{studentexercise}{quadratic variation after a nonlinear transform}
Let \(X_t=W_t^2-t\). Find \(\dd X_t\) and \([X]_t\). Is \(X\) a true
martingale on bounded horizons?
\end{studentexercise}
\begin{shortsolution}
\[
\dd X_t=2W_t\dd W_t,
\qquad
\boxed{[X]_t=4\int_0^tW_s^2\dd s.}
\]
Since \(\E\int_0^T4W_s^2\dd s=2T^2<\infty\), \(X\) is a square-integrable
martingale on every bounded horizon.
\end{shortsolution}

\subsection{Error diagnosis}

\begin{studentexercise}{find and correct four mistakes}
Suppose
\(\dd X_t=a_t\dd t+b_t\dd W_t\),
\(\dd W_t\dd B_t=\rho\dd t\),
\(\dd S_t/S_t=\mu\dd t+\sigma\dd W_t^{\Pp}\), and
\(\dd A_t/A_t=r\dd t\). Correct:
\[
\begin{aligned}
\dd f(t,X_t)
&=f_x\dd X_t+\frac12b_tf_{xx}\dd t,\\
\dd(W_tB_t)
&=W_t\dd B_t+B_t\dd W_t,\\
\text{zero drift}
&\Longrightarrow\text{true martingale},\\
\frac{\dd(S_t/A_t)}{S_t/A_t}
&=\sigma\dd W_t^{\Pp}.
\end{aligned}
\]
\end{studentexercise}
\begin{shortsolution}
\[
\begin{aligned}
\dd f(t,X_t)
&=\left(f_t+a_tf_x+\frac12b_t^2f_{xx}\right)\dd t
+b_tf_x\dd W_t,\\
\dd(W_tB_t)
&=W_t\dd B_t+B_t\dd W_t+\rho\dd t,\\
\text{zero drift}
&\Longrightarrow\text{local martingale first},\\
\frac{\dd(S_t/A_t)}{S_t/A_t}
&=(\mu-r)\dd t+\sigma\dd W_t^{\Pp}.
\end{aligned}
\]
The missing ingredients were, respectively: the time derivative and squared
diffusion coefficient; covariation; an integrability check; and the
physical excess-return drift.
\end{shortsolution}

\subsection{A compact end-of-lecture checklist}

\begin{keyidea}{Eight checks for every \(\Ito\) calculation}
\begin{enumerate}
  \item Identify the time and state variables.
  \item Compute \(f_t,f_x,f_{xx}\) before substituting the process.
  \item Write every driving SDE and covariation.
  \item Use \((\dd W)^2=\dd t\) and
        \(\dd W\dd B=\rho\dd t\).
  \item Retain every term of order \(\dd t\).
  \item Group the final answer into drift and diffusion.
  \item Check the domain, regularity, and integrability assumptions.
  \item In a financial argument, state the probability measure and explain
        the financing convention.
\end{enumerate}
\end{keyidea}

\begin{warningbox}{When the classical formula is not enough}
The formula in this lecture assumes \(f\in C^{1,2}\). It does not apply
directly at the kink of
\[
f(x)=|x|
\qquad\text{or}\qquad
f(s)=(s-K)^+.
\]
For nonsmooth convex functions, a local-time term appears in the generalized
\(\Ito\)--Tanaka formula. Before differentiating, locate singularities and
ask whether the process can reach them.
\end{warningbox}

\begin{keyidea}{Final synthesis}
\begin{itemize}
  \item Brownian quadratic variation makes stochastic calculus second order.
  \item Covariation turns correlation into a \(\dd t\)-term.
  \item The \(\Ito\) formula is a chain rule plus a curvature correction.
  \item Logarithms and integrating factors solve GBM, OU, and moment problems
        for CIR.
  \item Products, ratios, discounting, realized variance, and delta hedging
        all use the same multiplication table.
  \item Zero drift implies a local martingale first; suitable integrability
        upgrades it to a true martingale.
\end{itemize}
\end{keyidea}

\end{document}
